The total number of elements in a group is called the \emph{order} of the 
group.  It is readily seen that the order of a subgroup divides the order of 
the whole
\SourcePage{436}{439}
group.  To prove this, consider a subgroup $\boldsymbol{H}$ of a group 
$\boldsymbol{G}$, and let $G_1$ be an element of $\boldsymbol{G}$ that does not 
belong to $\boldsymbol{H}$.  Multiplication of every element of 
$\boldsymbol{H}$ by $G_1$---on the right, for example---gives a set (or 
\emph{coset}) of elements denoted by $\boldsymbol{H}G_1$.  Every element of 
this coset plainly belongs to $\boldsymbol{G}$, but none belongs to 
$\boldsymbol{H}$.  Indeed, if two elements $H_a,H_b$ of $\boldsymbol{H}$ 
satisfied $H_aG_1=H_b$, it would follow that 
$G_1=H_a^{-1}H_b$; this would put $G_1$ in $\boldsymbol{H}$, contrary to our 
assumption.  In the same way, if $G_2$ belongs to $\boldsymbol{G}$ but to 
neither $\boldsymbol{H}$ nor $\boldsymbol{H}G_1$, then none of the elements of 
$\boldsymbol{H}G_2$ belongs to either of those two sets.  Continuing in this 
fashion eventually exhausts all elements of the finite group 
$\boldsymbol{G}$.  Thus its elements are partitioned into sets (the 
\emph{cosets} of $\boldsymbol{H}$ in $\boldsymbol{G}$)
\[
\boldsymbol{H},\ \boldsymbol{H}G_1,\ \boldsymbol{H}G_2,\ldots,\boldsymbol{H}G_m,
\]
each containing $h$ elements, where $h$ is the order of the subgroup 
$\boldsymbol{H}$.  Consequently the order of $\boldsymbol{G}$ is $g=hm$, 
which proves the assertion.  The integer $m=g/h$ is called the \emph{index} 
of $\boldsymbol{H}$ in $\boldsymbol{G}$.
